% TOP_PROBLEM: {"id":30007572,"problem_number":"LOCAL-30007572","title":"Robust dynamic mechanisms","category_id":21,"category":{"id":21,"name":"economics","display_name":"Economics","description":"Open economic theory statements, published research agendas, and proposed scoped specifications; question provenance and historical priority are qualified per record.","slug":"economics","order_index":21,"created_at":"2026-10-08T00:00:00Z"},"status":"research_question","external_url":null,"legacy_ids":[],"published":false,"statement":"Characterize robust dynamic truthfulness and revenue guarantees when mechanisms must work for an ambiguity set of type-transition laws instead of one shared prior.","economics":{"original_id":61,"field":"Mechanism design and market design","kind":"theoretical","formalization_basis":"adapted_research_question","source_status":"research_agenda","priority_status":"earliest_found","priority_note":"This is the earliest source in which the relevant agenda was verified during this audit; absolute priority has not been established.","earliest_verified_reference":{"authors":["Dirk Bergemann","Juuso Välimäki"],"title":"Dynamic Mechanism Design: An Introduction","year":2018,"url":"https://cowles.yale.edu/sites/default/files/2022-09/d2102-r.pdf","doi":"10.1257/jel.20180892","locator":"§8 Concluding Remarks, pp.49–51 of June19,2018 revision","evidence_summary":"The conclusion explicitly identifies weaker common-prior assumptions, sequential participation, outside resale, and limited adoption of long-term contracts as research questions."},"current_reference":{"authors":["Dirk Bergemann","Juuso Välimäki"],"title":"Dynamic Mechanism Design: An Introduction","year":2018,"url":"https://cowles.yale.edu/sites/default/files/2022-09/d2102-r.pdf","doi":"10.1257/jel.20180892","locator":"§8 Concluding Remarks, pp.49–51 of June19,2018 revision","evidence_summary":"The conclusion explicitly identifies weaker common-prior assumptions, sequential participation, outside resale, and limited adoption of long-term contracts as research questions."},"formalization_note":"The ambiguity-set optimization is an adaptation. The paper does not state this exact ratio frontier; first-version priority before the read2018 revision is not established."},"statusQualification":"Published question or research agenda verified at the cited locator. The mathematical specification is adapted for this catalogue and includes additional assumptions; source publication does not establish first-ever priority or certify this exact model remains unsolved. The ambiguity-set optimization is an adaptation. The paper does not state this exact ratio frontier; first-version priority before the read2018 revision is not established.","statusReviewedAt":"2026-10-08"}
% FIRST_POSTED: 2026-10-08
% ENTRY_KIND: statement_only
% !TeX program = lualatex
\documentclass[11pt]{article}
\usepackage{fontspec}
\usepackage[a4paper,margin=25mm]{geometry}
\usepackage{amsmath,amssymb,mathtools}
\usepackage{xurl}
\usepackage[hidelinks]{hyperref}
\newcommand{\sref}[2]{\hyperlink{source:#1:#2}{[S#2]}}
\setlength{\emergencystretch}{3em}
\begin{document}
% problemId:30007572
\section*{Robust dynamic mechanisms}\label{30007572}
\subsection{Definitions and mathematical statement}
Consider a finite-horizon allocation problem with agents \(i=1,\ldots,n\), periods \(t=1,\ldots,T\), finite private type spaces \(\Theta_i\), feasible allocation set \(X\), and quasilinear discounted utility \(\sum_t\beta^{t-1}[v_i(x_t,\theta_{it})-p_{it}]\), where \(\beta\in(0,1]\). The designer knows an ambiguity set \(\mathcal P\) of joint type-transition laws, rather than one shared prior. A direct mechanism maps public report histories into allocations and transfers. Require truthful reporting to maximize every agent's expected continuation utility at every private history under every law \(P\in\mathcal P\), and require voluntary continuation under those same laws. Let \(\mathrm{Rev}_P(M)\) be expected discounted transfers and \(\mathrm{OPT}_P\) the revenue of the optimal dynamic mechanism when \(P\) is known. For a class where \(\mathrm{OPT}_P>0\), define
\[
 r^*(\mathcal P)=\sup_{M\text{ feasible and robustly truthful}}
 \inf_{P\in\mathcal P}\frac{\mathrm{Rev}_P(M)}{\mathrm{OPT}_P}.
\]
Characterize this guarantee and implementable allocation formats for ambiguity sets defined by known marginal distributions, transition bounds, or moments. Identify which restrictions permit positive guarantees and when heterogeneity of beliefs destroys the gains from intertemporal screening. Specify how zero-probability histories are handled and what information is available to each participant. Existing prior-independent dynamic mechanisms supply special cases; the target is an explicit robustness frontier, adapted from the survey’s weaker-information agenda.

\par\textbf{Formulation and scope.} The ambiguity-set optimization is an adaptation. The paper does not state this exact ratio frontier; first-version priority before the read2018 revision is not established.

\par\textbf{Open-question provenance.} An explicit question or research agenda was verified at the source locator. This does not by itself certify that every later version remains unresolved \sref{30007572}{1}.

\par\textbf{Historical priority.} This is the earliest source in which the relevant agenda was verified during this audit; absolute priority has not been established.

\subsection{Short English statement}
Characterize robust dynamic truthfulness and revenue guarantees when mechanisms must work for an ambiguity set of type-transition laws instead of one shared prior.

\subsection{Sources}
\begin{itemize}\raggedright
\item[\textnormal{[S1]}]\hypertarget{source:30007572:1}{} Dirk Bergemann, Juuso Välimäki. 2018. Dynamic Mechanism Design: An Introduction. §8 Concluding Remarks, pp.49–51 of June19,2018 revision.
\url{https://cowles.yale.edu/sites/default/files/2022-09/d2102-r.pdf}.
DOI: 10.1257/jel.20180892.
{\small\textit{Earliest verified question/agenda reference; historical priority qualified below. The conclusion explicitly identifies weaker common-prior assumptions, sequential participation, outside resale, and limited adoption of long-term contracts as research questions.}}
\end{itemize}
\end{document}
